Many differential equations begin with language such as "the rate of change is proportional to..." or "the rate depends on the current amount." A clean modelling workflow helps turn those words into mathematics.

\begin{enumerate}[leftmargin=*]
  \item \textbf{Define the variable.} State what $y$ represents and what the independent variable is. For example, $y$ could be a population and $x$ could be time in years.
  \item \textbf{Translate the rate statement.} If the growth rate is proportional to the amount present, write
  \[
  \frac{dy}{dx}=ky.
  \]
  If the rate depends on both the amount present and remaining capacity, write a logistic-style equation.
  \item \textbf{Solve the differential equation.} Use separation of variables or a known standard form.
  \item \textbf{Apply the condition.} Use initial data such as $y(0)=y_0$ to determine the constant.
  \item \textbf{Interpret constants.} Explain what $k$, $r$, or $M$ mean in the situation.
  \item \textbf{Check the answer against the story.} Ask whether the function stays positive when it should, whether units are sensible, and whether the long-run behavior is realistic.
\end{enumerate}

\subsubsection*{Useful sentence patterns}
\begin{itemize}[leftmargin=*]
  \item "Rate proportional to the current amount" usually leads to $\dfrac{dy}{dx}=ky$.
  \item "Rate proportional to the difference from a limiting value" often leads to $\dfrac{dy}{dx}=k(M-y)$.
  \item "Growth slows as the population approaches a maximum level" suggests a logistic model.
\end{itemize}

\subsubsection*{Mechanics substitution choices}
When a modelling question moves into Extension~2 mechanics, the key choice is often how to rewrite acceleration.

\begin{itemize}[leftmargin=*]
  \item Use
  \[
  a=\frac{dv}{dt}
  \]
  when the equation is naturally expressed in terms of time.
  \item Use
  \[
  a=v\frac{dv}{dx}
  \]
  when the question connects velocity and displacement directly.
  \item Use
  \[
  a=\frac{d}{dx}\left(\frac{1}{2}v^2\right)
  \]
  when an energy-style integral is the cleanest route.
\end{itemize}

These choices appear repeatedly in resistive motion, gravitation, and spring-motion models.
